\documentclass[a4paper,12pt]{book}
\usepackage[T1]{fontenc}
\usepackage[rm=lining,sf=lining,tt=lining]{cfr-lm}
\usepackage{silence}
\WarningFilter{nfssext-cfr}{Overwriting kernel definition of \swshape}
\WarningFilter{nfssext-cfr}{Overwriting the format's definition of \DeclareEncodingSubset}
\WarningFilter{nfssext-cfr}{Patching font initialisation macro for serif.}
\WarningFilter{nfssext-cfr}{Patching font initialisation macro for sans.}
\WarningFilter{nfssext-cfr}{Patching font initialisation macro for typewriter.}
\WarningFilter{spanish}{faked small caps}
\usepackage[utf8]{inputenc}
\usepackage[spanish,es-noquoting]{babel}

% Evita pedir "versalitas en negrita" (bx/sc), que no está disponible en Latin Modern.
\DeclareTextFontCommand{\textsc}{\scshape\mdseries}

\usepackage{fancyhdr}
\usepackage{amsmath, mathtools}
\usepackage{amssymb}
\usepackage{eurosym}    
\usepackage{graphicx}
\usepackage{float} 
\usepackage{vmargin}
\usepackage{cancel}
\usepackage{caption}
\captionsetup{justification=raggedright,singlelinecheck=false}
\usepackage{subcaption}
\usepackage{appendix}   
\usepackage[pdftex]{hyperref}   
\hypersetup{colorlinks=true,linkcolor=black,linkcolor=black,filecolor=black,urlcolor=black, citecolor=black}    \urlstyle{same} %hipervinculos en negro y sin subrayar
\usepackage[nottoc,notlot,notlof]{tocbibind}
\usepackage{listings}   
\usepackage{xcolor} 
\usepackage{wrapfig}    
\usepackage{transparent}
\usepackage{import}
\usepackage{verbatim}
\usepackage{afterpage}
\usepackage{emptypage}
\usepackage{etoolbox}
\usepackage{parskip}
\usepackage{multirow}
\usepackage{booktabs}
\usepackage{tabularx}
\usepackage{colortbl}
\usepackage{enumitem}
\usepackage{bm}
\usepackage{tikz}
\usepackage{pdfpages}
% TikZ shapes (p.ej., cylinder) usados en algunos diagramas
\usetikzlibrary{shapes.geometric, shapes.symbols, positioning, fit, backgrounds, shadows}

\raggedbottom                        % evita Underfull \vbox en páginas con poco contenido
\setlength{\emergencystretch}{3em}  % da margen extra a LaTeX para evitar Overfull \hbox borderline

\makeatletter
\definecolor{nar}{HTML}{FF8000}
\renewcommand\mainmatter{
    \clearpage
  \@mainmattertrue
  \pagenumbering{arabic}}
% Igual que \mainmatter: sin \cleardoublepage, para no dejar una página en
% blanco tras la bibliografía.
\renewcommand\backmatter{
    \clearpage
  \@mainmatterfalse}

\renewcommand\lstlistingname{Programación}
\renewcommand\lstlistlistingname{Índice de programación}
\addto\captionsspanish{\renewcommand*\contentsname{Índice}}
\addto\captionsspanish{\renewcommand{\tablename}{Tabla}  
          \renewcommand{\listtablename}{Índice de tablas}}

% cambiar márgenes chapters
\patchcmd{\@makechapterhead}{50\p@}{\chapheadtopskip}{}{}% Space from top of page to CHAPTER X
\patchcmd{\@makechapterhead}{20\p@}{\chapheadsep}{}{}% Space between CHAPTER X and CHAPTER TITLE
\patchcmd{\@makechapterhead}{40\p@}{\chapheadbelowskip}{}{}% Space between CHAPTER TITLE and text
% --- Patch \chapter*
\patchcmd{\@makeschapterhead}{50\p@}{\chapheadtopskip}{}{}% Space from top of page to CHAPTER TITLE
\patchcmd{\@makeschapterhead}{40\p@}{\chapheadbelowskip}{}{}% SPace between CHAPTER TITLE and text
\makeatother
% Set new lengths
\newlength{\chapheadtopskip}\setlength{\chapheadtopskip}{20pt}
\newlength{\chapheadsep}\setlength{\chapheadsep}{40pt}
\newlength{\chapheadbelowskip}\setlength{\chapheadbelowskip}{15pt}

\renewcommand{\arraystretch}{1.3}

%Dimensiones
\textwidth = 16cm
\textheight = 25cm
\topmargin = 1cm
\oddsidemargin = 2.5cm
\evensidemargin = 2.5cm
\headsep = 1cm
\footskip = 1.1cm
\setlength{\headheight}{32pt}

%Encabezado y pie de página

\makeatletter
\renewcommand\chapter{\if@openright\cleardoublepage\else\clearpage\fi
\thispagestyle{fancy}% original style: plain
\global\@topnum\z@
\@afterindentfalse
\secdef\@chapter\@schapter}
%la siguiente línea es para que no haya páginas en blanco entre los capítulos
\patchcmd{\chapter}{\if@openright\cleardoublepage\else\clearpage\fi}{}{}{}
%para quitar lo de "capítulo 1"
\def\@makechapterhead#1{%
  \vspace*{50\p@}%
  {\parindent \z@ \raggedright \normalfont
    \ifnum \c@secnumdepth >\m@ne
      \if@mainmatter
        %\huge\bfseries \@chapapp\space \thechapter
        \Huge\bfseries \thechapter.\space%
        %\par\nobreak
        %\vskip 20\p@
      \fi
    \fi
    \interlinepenalty\@M
    \Huge \bfseries #1\par\nobreak
    \vskip 40\p@
  }}
\makeatother %para usar lo mismo en la primera pag del capitulo


\pagestyle{fancy}
\fancyhf{}
\fancyhead[LE,RO]{\transparent{0.6}\includegraphics[width=3cm]{portada/logosportada/logoetsit.png}}
\fancyhead[RE,LO]{\transparent{0.6}\thepage}
\chead{\transparent{0.6}Josué Bohórquez Cristóbal}
\fancyfoot[C]{\scriptsize Diseño e implementación del motor cinemático y navegación de un robot hexápodo en ROS~2}
\renewcommand{\headrulewidth}{0.4pt} % grosor de la línea de la cabecera
\renewcommand{\footrulewidth}{0.4pt} % grosor de la línea del pie

\begin{document}

%Numeracion romana
\frontmatter

%Portada
\begin{titlepage}
\pagecolor{nar}
\afterpage{\nopagecolor}
    \begin{center}
        \vspace{1cm}
        \LARGE{\textsc{Universidad Politécnica de Madrid}}\par
        \vspace{0.3cm}
        \large{\textsc{Escuela Técnica Superior de Ingenieros de Telecomunicación}} \par
    \end{center}
    
    \begin{figure}[H]
        \vspace{1cm}
        \centering
        \includegraphics[width=10cm]{portada/logosportada/byn.png}
    \end{figure}

    \begin{center}
        \vspace{1.5cm}
        \Large{\textsc{Grado en Ingeniería de Tecnologías y Servicios de
Telecomunicación}}\par
        \vspace{0.2cm}
        \LARGE{\textsc{Trabajo Fin de Grado}}\par
        \vspace{2.5cm}

        \Huge{\textsc{Diseño e implementación del motor cinemático y navegación de un robot hexápodo en 
ROS 2 }}

        \vspace*{2.5cm}
                  \large{\textsc{ Josué Bohórquez Cristóbal }}
        \par\vspace{0.5cm}
                  \normalsize{\textsc{Tutor: D. Álvaro Gutiérrez Martín}}
        \par\vspace{2cm}
               \small{\the\year}
    \end{center}
\end{titlepage}
\empty

% Segunda portada
\begin{titlepage}
    \begin{center}
        \vspace{1cm}
        \LARGE{\textsc{Universidad Politécnica de Madrid}}\par
        \vspace{0.3cm}
        \large{\textsc{Escuela Técnica Superior de Ingenieros de Telecomunicación}} \par
    \end{center}
    
    \begin{figure}[H]
        \vspace{1cm}
        \centering
        \includegraphics[width=10cm]{portada/logosportada/logoetsit.png}
    \end{figure}

    \begin{center}
        \vspace{1.5cm}
        \Large{\textsc{Grado en Ingeniería de Tecnologías y Servicios de Telecomunicación}}\par
        \vspace{0.2cm}
        \LARGE{\textsc{Trabajo Fin de Grado}}\par
        \vspace{2.5cm}
        \Huge{\textbf{Diseño e implementación del motor cinemático y navegación de un robot hexápodo en 
ROS 2 }}

        \vspace*{2.5cm}
                  \large{ Josué Bohórquez Cristóbal }
        \par\vspace{0.5cm}
                  \normalsize{Tutor: D. Álvaro Gutiérrez Martín}
        \par\vspace{2cm}
               \small{\the\year}
    \end{center}
\end{titlepage}
\empty

% Resumen y agradecimientos 
\chapter*{}
\addcontentsline{toc}{chapter}{Resumen y Palabras Clave}
% Resumen en español
\Large{\textsc{\textcolor{nar}{Resumen}}} \par
\vspace{0.2cm}
\normalsize{Número máximo de palabras: 500

...}

\vspace{0.9cm}

% Summary (resumen en inglés)
\Large{\textsc{\textcolor{nar}{Summary}}} \par
\vspace{0.2cm}
\normalsize{Maximum number of words: 500

...}
\vspace{0.9cm}

% Palabras clave
\Large{\textsc{\textcolor{nar}{Palabras Clave}}} \par \par
\vspace{0.2cm}
\normalsize{Deben reflejar el contenido del trabajo, deberían servir para localizar el TFG mediante búsqueda bibliográfica...} \par
\vspace{0.8cm}

\Large{\textsc{\textcolor{nar}{Keywords}}} \par
\vspace{0.2cm}
\normalsize{...} \par

\newpage
\chapter{Agradecimientos}
...

% Indice, lista de figuras y de tablas
\newpage
\begin{center}
   \tableofcontents
\end{center}
\newpage
\begin{center}
   \listoffigures
   \listoftables
\end{center}
\newpage

% Lista de acrónimos
    \chapter{Lista de acrónimos}

\begin{description}
    \item[API] \textit{Application Programming Interface}.
    \item[CCW] \textit{Counterclockwise}.
    \item[CPG] \textit{Central Pattern Generator}.
    \item[CW] \textit{Clockwise}.
    \item[DC] \textit{Direct Current}.
    \item[DDS] \textit{Data Distribution Service}.
    \item[EEPROM] \textit{Electrically Erasable Programmable Read-Only Memory}.
    \item[ESP-IDF] \textit{Espressif IoT Development Framework}.
    \item[GPIO] \textit{General Purpose Input/Output}.
    \item[IK] \textit{Inverse Kinematics}.
    \item[IMU] \textit{Inertial Measurement Unit}.
    \item[LiDAR] \textit{Light Detection and Ranging}.
    \item[LiPo] \textit{Lithium Polymer}.
    \item[NUKE] \textit{Nearly Universal Kinematics Engine}.
    \item[PCB] \textit{Printed Circuit Board}.
    \item[PWM] \textit{Pulse Width Modulation}.
    \item[RAM] \textit{Random Access Memory}.
    \item[ROS] \textit{Robot Operating System}.
    \item[TFG] Trabajo Fin de Grado.
    \item[TTL] \textit{Transistor-Transistor Logic}.
    \item[UART] \textit{Universal Asynchronous Receiver-Transmitter}.
\end{description}

\newpage

\mainmatter

\chapter{Introducción y objetivos}
    \import{capitulos/cap0/}{intro.tex}
\newpage

\chapter{Estado del arte}
    \import{capitulos/cap1/}{estadoDelArte.tex}
\newpage

\chapter{Marco Teórico y Fundamentos}
    \import{capitulos/cap2/}{fundamentos.tex}
\newpage

\chapter{Plataforma hardware e implementación}
    \import{capitulos/cap3/}{plataformahardware.tex}
\newpage

\chapter{Validación, conclusiones y líneas futuras}
    \import{capitulos/cap4/}{ejemploResultados.tex}
    \import{capitulos/cap5/}{conclusiones.tex}
    \import{capitulos/cap5/}{lineasfuturas.tex}

\newpage
\bibliographystyle{unsrt}
\bibliography{biblio}

\backmatter
\chapter{Anexo A: Aspectos éticos, económicos, sociales y ambientales}
    \import{capitulos/anexos/}{anexoA.tex}


\newpage
\chapter{Anexo B: Presupuesto económico}
    \import{capitulos/anexos/}{anexoB.tex}
\newpage

\chapter{Anexo C: Documentación técnica de la PCB principal}
\label{anexo:documentacion_pcb}
    \import{capitulos/anexos/}{anexoC.tex}
\newpage

\chapter{Anexo D: Manual de puesta en marcha del sistema}
\label{anexo:manual_firmware}
    \import{capitulos/anexos/}{anexoD.tex}
\newpage

\chapter{Anexo E: Capturas de la interfaz web}
\label{anexo:capturas_panel_web}
    \import{capitulos/anexos/}{anexoE.tex}
\newpage



\end{document}
